\section{The Personal Equation and Belief}

I often receive correspondence from readers of various quality. In a recent one the interlocutor assumed that my ``emotional attachment" to a tradition was preventing me from achieving detachment. Careful readers will not find that anywhere on Gornahoor. However, I noted that Bard — Google's new AI system — made this evaluation:

\begin{quotex}
Salvo's writing is often dense and difficult to understand, but he has a large following of readers who appreciate his unique perspective on the world.

\end{quotex}
Well, if such a high machine intelligence finds us difficult to understand, then a merely human intelligence can be forgiven for its ineptitude. So consider what follows to be a brief summary of what is being asserted in many and various ways.

\paragraph{Exoteric and Esoteric}
Living traditions have an exoteric and an esoteric teaching. It is a common misunderstanding to confound the two. The exoteric both hides and protects the esoteric teaching. Moreover, outside observers often see the egregore of the exoteric side, which they presume to be the true teaching. That is unfortunate, yet is understandable.

\paragraph{Metaphysics}
Metaphysics is as objective and knowable as any science, including physics. However, physics only concerns itself with the working or the physical world, or gross manifestation. Metaphysics, on the other hand, recognizes that reality is much vaster than what science can conceive. There are subtle states which depend on consciousness and can be explored by the intrepid inner researcher. Beyond that, there are many degrees of existence with various qualities. Poets, mystics, and saints — of many traditions — have described those states. Physics and metaphysics are not in conflict unless they intrude into the other's problem domain.

In this sense, Vladimir Solovyov asserted that Christianity did not introduce a new metaphysics. Rather it is base on Neoplatonic and Hermetic philosophy, which are the Western version of the Ancient Tradition. Its contribution comprises the historical events that culminated in the birth, death, and resurrection of Jesus.

\paragraph{History and Faith}
Although history is concerned with the gross world, unlike physics it cannot be known simply from external observations. The meaning of history requires an understanding of the intentions of the various actors.

Clearly, at that time some believed that Jesus was the Messiah, others that he was an enemy of the state, but most just saw an itinerant preacher. So which version are we to accept? Thomas Aquinas explains:

\begin{quotex}
Believing is an act of the intellect assenting to the divine truth by command of the will.

\end{quotex}
There is nothing there about emotions or irrationality. Of course, exoteric believers will sound very emotional and sentimental about their beliefs, but not those with a deeper understanding. There are two aspects to Thomas's insight.

First of all, belief is a deliberate choice, i.e., an act of will. That is not at all an emotional reaction. Nevertheless, emotions may follow. However, these are higher emotions that follow from the intellect; they do not control the choice as the lower emotions do.

Second, the intellect assents implying that the belief is not compellingly false nor implausible. Without going into complete detail, I can assert this. The intellect sees that mankind is fallen in some sense from its primordial glory. We can also see that there is not worldly solution: no science, pharmaceutical, eugenic, psychological, economic, or social solution can ever remedy that. Only the insertion of a force from outside the system can do so. That may not be a proof for you, but it is a justified belief.

\paragraph{Conclusion}
So I willingly embrace an exoteric tradition that is helpful to me personally and fits my personal equation. It has spiritually nourished my direct ancestors for 1500 years, so I am loathe to reject it arbitrarily. As I wrote above, there is simply no need to. My reader claims that, ``I'm on the side of Man's increasing consciousness, not specifically the Christians, the Jews, the Hindus etc."

That makes as much sense as going to a restaurant and telling the server: ``I am in favor of eating, but not specifically Italian, German, Japanese; just give me food."

Just as healthy food nourishes the body, whatever type of cuisine it is, so does a healthy tradition nourish the spirit. For me to abandon my Tradition to chase after some vague ``increasing consciousness" means that I would have to reject the metaphysical teachings of Augustine, Boethius, Bernard, Dante, Thomas Aquinas, Eckhart, and so many others. But ironically, they are the very teachers of ``increasing consciousness."

Moreover, there is also a rich source of teachings on spiritual practice, prayer, meditation, contemplation, and so on. It is hubris for anyone to think he or she can recreate all that de novo.



\flrightit{Posted on 2023-04-04 by Cologero }

\begin{center}* * *\end{center}

\begin{footnotesize}\begin{sffamily}



\texttt{taivasvaeltaja on 2023-04-09 at 16:38 said: }

Even though it may sound sentimental, the gift of tears is still a password for higher realms for many. When Jesus wept, did hs tears make oceans?


\hfill

\texttt{Sylvan Savant on 2023-04-09 at 18:23 said: }

If there's no orthodoxy, how could there be any heresy? If there is no right path to the destination, then how could one go astray? 

Excellent quote by Thomas Aquinas. I find it very helpful in my current situation, thanks.


\hfill

\texttt{Taivasvaeltaja on 2023-10-23 at 09:28 said: }

``In this sense, Vladimir Solovyov asserted that Christianity did not introduce a new metaphysics. Rather it is base on Neoplatonic and Hermetic philosophy, which are the Western version of the Ancient Tradition. Its contribution comprises the historical events that culminated in the birth, death, and resurrection of Jesus."

I've become to think that perhaps Christianity sort of restricted its metaphysical teaching to a certain level of understanding in order to make men focus on the process of salvation, where all the theology of the churches then sprang out.

That which does not make me doubt in any sense is that the symbolical and mythical life of Jesus Christ was a momentous happening on planet earth and in which the symbols that define the Christic Mythos, the turning of the wheel of involution towards the ascending arch of spiritual evolution, liberation of maya and matter as the bodies of Christ were crucified while his Solar Nature laughs upon the Cross. As Christ is the Light and Lucifer the Light-Bearer, the mental body of Christ is Sophia-Lucifer – the one whom even Gnostics struggled to understand – and the bodily cross of man is Lucifer's prison. Only the Logos could have done the turning of the wheel of involution, because ``The Father and I are One."


\end{sffamily}\end{footnotesize}
